% instancias de representacion
\thispagestyle{empty}
\PageDecor{4}
\SectionTag{SECCIÓN INFORMATIVA}
\DocHeading{Instancias de Representación Estudiantil}
\DocSubheading{Las y los estudiantes cuentan con distintas instancias de representación:}
\DocItem{\textbf{Consejo Directivo:} es el órgano supremo de gobierno de la institución. Está compuesto por cinco consejeros por el claustro docente, cinco por el estudiantil, uno por el no docente y dos por el graduado. Es presidido por la Rectora de la institución. Cada dos años se celebran elecciones a Consejo Directivo.}
\DocItem{\textbf{Centro de Estudiantes (CD):} es una organización conformada por y para los y las estudiantes. Es legitimado por el voto de sus compañeras y compañeros, posee objetivos, actividades y propuestas, consta de una estructura organizacional, con derechos y responsabilidades. Lo componen la presidencia, la vice-presidencia, las distintas secretarías y vocalías. Su organización está regida por el Estatuto del Centro de Estudiantes, aprobado al momento de su creación.}
\DocItem{\textbf{Asamblea Estudiantil:} es un espacio de debate y decisión colectiva, convocada anualmente o extraordinaria según su necesidad, en la cual cada estudiante cuenta con voz y voto sobre los temas que se traten. El modo de convocatoria y su esquema está estipulado también en el Estatuto del Centro de Estudiantes. Lo resuelto en Asamblea por la comunidad estudiantil afecta a todas las demás instancias de representación estudiantil (Centro, Consejeros y Delegados/as).}
\DocItem{\textbf{Cuerpo de Delegadas y Delegados:} su función principal es representar al año y división al que pertenezcan. Otras de sus funciones son informar a sus respectivos cursos las resoluciones y medidas que adopte el CE y demás situaciones de la vida académica, participar con voz, pero sin voto de las reuniones y presentar propuestas, inquietudes, etcétera. Se conforma por un representante titular y otro suplente de cada año y división. Cada delegada o delegado es elegido por voto simple y directo dentro de los primeros 30 días del ciclo lectivo.}
\DocItem{\textbf{Protocolo de Genero:} espacio de intervención institucional ante situaciones de violencia de género, acoso sexual y discriminación de género. El mismo está conformado por un estudiante, docentes y no docentes del profesorado.}
